\chapter{Advanced Features}
\label{ch:advanced}

This chapter covers advanced language features that build upon the core language.

\section{Generics}

Generics allow functions and types to be parameterized over types.

\subsection{Generic Functions}

\begin{lstlisting}
fn identity<T>(value: T) -> T {
    return value
}

fn swap<T, U>(pair: (T, U)) -> (U, T) {
    return (pair.1, pair.0)
}

fn first<T>(items: []T) -> T? {
    if array.len(items) == 0 {
        return nil
    }
    return items[0]
}

// Usage - type inferred
identity(42)           // T = int
identity("hello")      // T = string
swap((1, "one"))       // T = int, U = string
\end{lstlisting}

\subsection{Generic Types}

\begin{lstlisting}
struct Pair<T, U> {
    first: T
    second: U
}

struct Stack<T> {
    items: []T
}

fn (s: *Stack<T>) push<T>(item: T) {
    s.items = array.push(s.items, item)
}

fn (s: *Stack<T>) pop<T>() -> T? {
    if array.len(s.items) == 0 {
        return nil
    }
    s.items, item = array.pop(s.items)
    return item
}

// Usage
stack: Stack<int> = Stack{items: []}
stack.push(1)
stack.push(2)
value = stack.pop()  // 2
\end{lstlisting}

\subsection{Type Constraints}

Constrain generic types to those implementing certain behaviors:

\begin{lstlisting}
constraint Comparable {
    fn compare(other: Self) -> int
}

constraint Printable {
    fn to_string() -> string
}

constraint Hashable {
    fn hash() -> int
}

// Constrained generic
fn max<T: Comparable>(a: T, b: T) -> T {
    if a.compare(b) > 0 {
        return a
    }
    return b
}

// Multiple constraints
fn process<T: Comparable + Printable>(items: []T) {
    sorted = array.sort(items)
    for item in sorted {
        io.println(item.to_string())
    }
}
\end{lstlisting}

\subsection{Where Clauses}

For complex constraints:

\begin{lstlisting}
fn merge<K, V>(a: [K:V], b: [K:V]) -> [K:V]
    where K: Hashable + Comparable {
    result = a
    for k, v in b {
        result[k] = v
    }
    return result
}
\end{lstlisting}

\section{Constraints (Traits)}

Constraints define shared behavior that types can implement.

\subsection{Defining Constraints}

\begin{lstlisting}
constraint Iterator<T> {
    fn next() -> T?
    fn has_next() -> bool
}

constraint Serializable {
    fn serialize() -> []byte
    fn deserialize(data: []byte) -> (Self, Error?)
}

constraint Ordered {
    fn less_than(other: Self) -> bool
    fn greater_than(other: Self) -> bool
    fn equals(other: Self) -> bool
}
\end{lstlisting}

\subsection{Implementing Constraints}

\begin{lstlisting}
struct Counter {
    current: int
    max: int
}

impl Iterator<int> for Counter {
    fn next() -> int? {
        if self.current >= self.max {
            return nil
        }
        value = self.current
        self.current += 1
        return value
    }

    fn has_next() -> bool {
        return self.current < self.max
    }
}

// Usage
counter = Counter{current: 0, max: 5}
while counter.has_next() {
    io.println(counter.next())
}
\end{lstlisting}

\subsection{Built-in Constraints}

\begin{lstlisting}
constraint Eq {
    fn equals(other: Self) -> bool
}

constraint Ord: Eq {
    fn compare(other: Self) -> int  // -1, 0, or 1
}

constraint Hash {
    fn hash() -> int
}

constraint Clone {
    fn clone() -> Self
}

constraint Default {
    fn default() -> Self
}

constraint Display {
    fn display() -> string
}

constraint Debug {
    fn debug() -> string
}
\end{lstlisting}

\section{Union Types}

Union types represent values that can be one of several types.

\subsection{Simple Unions}

\begin{lstlisting}
type StringOrInt = string | int

fn process(value: StringOrInt) {
    match value {
        s: string => io.println("String: " + s)
        n: int => io.println("Int: " + to_string(n))
    }
}

process("hello")
process(42)
\end{lstlisting}

\subsection{Tagged Unions (Enums with Data)}

\begin{lstlisting}
enum Result<T, E> {
    Ok(T),
    Err(E)
}

enum Option<T> {
    Some(T),
    None
}

enum Json {
    Null,
    Bool(bool),
    Number(float),
    String(string),
    Array([]Json),
    Object([string:Json])
}

// Pattern matching
fn process_json(json: Json) {
    match json {
        Json.Null => io.println("null")
        Json.Bool(b) => io.println("bool: " + to_string(b))
        Json.Number(n) => io.println("number: " + to_string(n))
        Json.String(s) => io.println("string: " + s)
        Json.Array(items) => {
            for item in items {
                process_json(item)
            }
        }
        Json.Object(obj) => {
            for key, value in obj {
                io.print(key + ": ")
                process_json(value)
            }
        }
    }
}
\end{lstlisting}

\section{Pattern Matching}

Advanced pattern matching features.

\subsection{Destructuring}

\begin{lstlisting}
// Tuple destructuring
(x, y, z) = get_coordinates()

// Struct destructuring
User{name, email, age} = get_user()

// Nested destructuring
Point{x: Point{x: inner_x, y: _}, y: _} = nested_point

// Array destructuring
[first, second, ...rest] = items
[head, ...tail] = list
\end{lstlisting}

\subsection{Pattern Guards}

\begin{lstlisting}
match value {
    n if n < 0 => io.println("negative")
    n if n == 0 => io.println("zero")
    n if n > 0 and n < 100 => io.println("small positive")
    n if n >= 100 => io.println("large")
}
\end{lstlisting}

\subsection{Or Patterns}

\begin{lstlisting}
match status_code {
    200 | 201 | 204 => io.println("success")
    400 | 404 | 422 => io.println("client error")
    500 | 502 | 503 => io.println("server error")
    _ => io.println("unknown")
}
\end{lstlisting}

\subsection{Binding Patterns}

\begin{lstlisting}
match user {
    User{name: n, age: a} if a >= 18 => {
        io.println(n + " is an adult")
    }
    User{name: n, age: a} => {
        io.println(n + " is " + to_string(a) + " years old")
    }
}
\end{lstlisting}

\section{Operator Overloading}

Types can implement operators through constraints.

\begin{lstlisting}
constraint Add<T> {
    fn add(other: T) -> T
}

constraint Sub<T> {
    fn sub(other: T) -> T
}

constraint Mul<T> {
    fn mul(other: T) -> T
}

constraint Div<T> {
    fn div(other: T) -> T
}

struct Vector2 {
    x: float
    y: float
}

impl Add<Vector2> for Vector2 {
    fn add(other: Vector2) -> Vector2 {
        return Vector2{
            x: self.x + other.x,
            y: self.y + other.y
        }
    }
}

// Usage: a + b calls a.add(b)
v1 = Vector2{x: 1.0, y: 2.0}
v2 = Vector2{x: 3.0, y: 4.0}
v3 = v1 + v2  // Vector2{x: 4.0, y: 6.0}
\end{lstlisting}

\section{Compile-Time Evaluation}

Constant expressions evaluated at compile time.

\begin{lstlisting}
const PI = 3.14159265358979
const TAU = PI * 2.0
const MAX_BUFFER_SIZE = 1024 * 1024

// Compile-time functions
const fn factorial(n: int) -> int {
    if n <= 1 {
        return 1
    }
    return n * factorial(n - 1)
}

const FACTORIAL_10 = factorial(10)  // Computed at compile time
\end{lstlisting}

\section{Attributes}

Attributes modify declaration behavior.

\begin{lstlisting}
// Deprecation
#[deprecated("Use new_function instead")]
fn old_function() { }

// Conditional compilation
#[cfg(os = "linux")]
fn linux_specific() { }

#[cfg(debug)]
fn debug_only() { }

// Inlining hints
#[inline]
fn small_function() { }

#[inline(never)]
fn large_function() { }

// Testing
#[test]
fn test_something() { }

// Serialization
#[derive(Serialize, Deserialize)]
struct Config {
    #[rename("server_port")]
    port: int

    #[skip]
    internal_state: InternalState
}
\end{lstlisting}

\section{Macros}

Compile-time code generation.

\begin{lstlisting}
// Simple macro
macro debug!(expr) {
    io.println(stringify!(expr) + " = " + to_string(expr))
}

debug!(x + y)  // Expands to: io.println("x + y" + " = " + to_string(x + y))

// Variadic macro
macro println!(format, ...args) {
    io.printf(format + "\n", args...)
}

println!("Hello, %s!", name)

// Assert macro
macro assert!(condition, message = "assertion failed") {
    if not condition {
        panic(message + " at " + file!() + ":" + line!())
    }
}
\end{lstlisting}

\section{Unsafe Code}

For low-level operations.

\begin{lstlisting}
unsafe {
    // Raw pointer operations
    ptr = raw_pointer(data)
    value = *ptr
    *ptr = new_value

    // FFI calls
    result = c_function(arg1, arg2)
}

// Unsafe function
unsafe fn direct_memory_access(addr: usize) -> u8 {
    ptr = addr as *u8
    return *ptr
}
\end{lstlisting}

\begin{warning}
Unsafe code bypasses \haira{}'s safety guarantees. Use only when necessary and document thoroughly.
\end{warning}

\section{Foreign Function Interface (FFI)}

Calling C code.

\begin{lstlisting}
// Declare external function
extern "C" {
    fn printf(format: *i8, ...) -> i32
    fn malloc(size: usize) -> *void
    fn free(ptr: *void)
}

// Use in unsafe block
unsafe {
    ptr = malloc(1024)
    // Use memory...
    free(ptr)
}

// Wrapper for safe interface
fn safe_alloc(size: int) -> []byte {
    unsafe {
        ptr = malloc(size as usize)
        return bytes_from_raw(ptr, size)
    }
}
\end{lstlisting}

\section{Reflection}

Runtime type information.

\begin{lstlisting}
import "reflect"

fn describe(value: any) {
    t = reflect.type_of(value)
    io.println("Type: " + t.name)
    io.println("Kind: " + t.kind)

    if t.kind == "struct" {
        for field in t.fields {
            io.println("  " + field.name + ": " + field.type.name)
        }
    }
}

struct User {
    name: string
    age: int
}

describe(User{name: "Alice", age: 30})
// Output:
// Type: User
// Kind: struct
//   name: string
//   age: int
\end{lstlisting}

\section{Advanced Features Grammar}

\begin{lstlisting}[language=EBNF]
generic_params = "<" type_param { "," type_param } ">" ;

type_param = identifier [ ":" constraint_list ] ;

constraint_list = identifier { "+" identifier } ;

where_clause = "where" where_item { "," where_item } ;

where_item = identifier ":" constraint_list ;

constraint_def = "constraint" identifier [ generic_params ] "{"
                 { fn_signature }
                 "}" ;

impl_block = "impl" [ generic_params ] identifier "for" type
             [ where_clause ] "{" { fn_def } "}" ;

union_type = type { "|" type } ;

attribute = "#[" identifier [ "(" attr_args ")" ] "]" ;
\end{lstlisting}
